\chapter{Architecture and Methodology}
\section{System allocation}
The intended node uses an ESP32-S3 and MPU6050. The laptop gateway owns transport reception, local storage, explanation, and PDF generation. The mobile/display layer reads the gateway API. Figure~\ref{fig:architecture} separates the intended wearable from implemented gateway modules.
\begin{figure}[htbp]
\centering
\begin{tikzpicture}[
  stage/.style={draw=black!60, rounded corners=2pt, align=center,
    text width=2.05cm, minimum height=10mm, inner sep=3pt, font=\scriptsize},
  planned/.style={stage, dashed, fill=black!2},
  flow/.style={-{Stealth[length=3mm,width=2mm]}, line width=.85pt,
    draw=black!70},
  plannedflow/.style={flow, dashed}
]
\node[font=\small\bfseries] at (0,1.65)
  {Wearable / MCU path --- intended; hardware validation pending};

\node[planned] (sensor) at (-5.8,.55) {MPU6050\\acquisition};
\node[planned] (buffer) at (-2.9,.55) {Timed window\\buffer};
\node[planned] (gate) at (0,.55) {Layer 1\\threshold gate};
\node[planned] (cnn) at (2.9,.55) {INT8 CNN};
\node[planned] (transport) at (5.8,.55) {Event\\transport};
\draw[plannedflow] (sensor.east) -- (buffer.west);
\draw[plannedflow] (buffer.east) -- (gate.west);
\draw[plannedflow] (gate.east) -- (cnn.west);
\draw[plannedflow] (cnn.east) -- (transport.west);

\node[font=\small\bfseries] at (0,-.45)
  {Laptop gateway path --- software modules available};
\node[stage] (parser) at (5.8,-1.4) {Event\\parser};
\node[stage] (storage) at (2.9,-1.4) {Local JSON\\storage};
\node[stage] (explain) at (0,-1.4) {Explanation\\and report};
\node[stage] (api) at (-2.9,-1.4) {Local REST\\API};
\node[stage] (client) at (-5.8,-1.4) {Display-only\\client};
\draw[flow] (parser.west) -- (storage.east);
\draw[flow] (storage.west) -- (explain.east);
\draw[flow] (explain.west) -- (api.east);
\draw[flow] (api.west) -- (client.east);
\draw[plannedflow] (transport.south) -- (parser.north);

\end{tikzpicture}
\caption{Conceptual system architecture and data flow. Dashed boxes and arrows mark intended, hardware-unvalidated stages; solid boxes and arrows mark available gateway software.}
\label{fig:architecture}
\end{figure}


\section{Sampling and input contract}
Let $f_s$ be sampling frequency and $N$ the number of samples. The nominal window duration is $T=N/f_s$. Existing preprocessing sets $f_s=100\,\mathrm{Hz}$ and $N=200$, hence $T=2\,\mathrm{s}$. Model arrays have six channels ordered as acceleration then angular-rate axes. The planned logger targets 200 Hz acquisition; it is not evidence that the model was trained at that rate.

A 200-sample stream at 200 Hz has one second of coverage. A three-second 200 Hz input requires 600 samples. Consequently, the historical combination ``200 Hz, three seconds, 200$\times$6'' is rejected. This report retains the executable model convention and requires a versioned resampling adapter for any higher-rate logger. It does not silently alter the trained model.

The source preparation retains sensor counts in the selected columns; it does not establish a calibrated physical-unit conversion. Before connecting firmware, reconcile accelerometer/gyro scaling, axis signs, normalization, interpolation, anti-aliasing, and quantization. Timestamp gaps, clipping, and non-finite input must be explicit rejection or quality flags.

\section{Gate and classifier}
The acceleration magnitude is
\[
 a_m=\sqrt{a_x^2+a_y^2+a_z^2}.
\]
Current host gate code triggers when magnitude is at least 2.5 g and the continuous above-threshold duration exceeds 300 ms. This differs from older prose describing a shorter-than-300 ms spike. The implemented behavior must be evaluated against annotated falls before either condition is accepted as the physical design. No pre-impact lead time is established.

The CNN architecture uses two Conv1D stages, pooling, global average pooling, a dense hidden stage, and two-class softmax output. Class 1 represents FALL and class 0 represents non-fall. For a frozen decision threshold $\tau$, the binary decision is
\[
\hat y=\begin{cases}
1 & \text{if }P(\mathrm{FALL})\ge\tau \quad\text{(fall)},\\
0 & \text{otherwise} \quad\text{(non-fall)}.
\end{cases}
\] Threshold tuning belongs on validation data; test labels must not select $\tau$.

Full-integer quantization maps a float input to $q=\operatorname{clip}(\operatorname{round}(x/s)+z)$ using model-specific scale $s$ and zero point $z$. The loaded artifact must supply these values; they cannot be borrowed across model generations. Quantized output dequantization precedes threshold comparison.

\section{Holdout and evaluation design}
The training source uses grouped holdouts and asserts subject-set disjointness. Group-based fractions need not yield identical fractions of windows, and the helper is not repeated cross-validation. The first source split reserves a test group fraction of 0.2; the second split uses a validation fraction of $0.1/0.9$ within the retained groups. Do not describe this code as a guaranteed 80/10/10 window split. The October notebook run must be traced separately from the CLI source.

Retain exact subject/episode IDs, class counts, seeds, preprocessing revision, and hashes for each split. Trial-level labels inherited by every window can include pre/post-fall periods; event-level ground truth must eventually resolve onset and contact rather than treating every such window as a separate fall.

\section{Event and explanation contract}
Wire Format v1 records event identity, device/version, uint32 device timestamp, $P(\mathrm{FALL})$, and six signed peak features in g and degrees/s. It omits the raw waveform. Therefore a v1 live event cannot alone drive the waveform CNN explanation. An approved extension or linked waveform record requires synchronized firmware and gateway changes.

For a supplied raw window, the implemented attribution has the form
\[
 A_{t,c}=X_{t,c}\frac{\partial P(\mathrm{FALL})}{\partial X_{t,c}},\qquad
 r_c=\frac{\sum_t |A_{t,c}|}{\sum_{t,c}|A_{t,c}|}.
\]
These normalized channel magnitudes are saliency summaries, not causal biomechanics or exact Shapley values. The fixed temporal segments in the code are not event-aligned impact annotations.
\section{Requirements and acceptance evidence}
These requirements describe intended behavior, not completed certification.
\begin{description}
\item[FR-01] The logger shall preserve timestamped six-channel recordings and sample-quality metadata.
\item[FR-02] The classifier shall use a versioned sampling, scaling, shape, and threshold contract tied to a model hash.
\item[FR-03] The event pipeline shall preserve unique event identity, confidence, units, and model/device provenance.
\item[FR-04] The gateway shall store events locally and expose reports to a display-only client.
\item[FR-05] Volunteer recording shall require both institutional approval and logger validation.
\item[NFR-01] The planned acquisition target shall be 200 Hz (5 ms nominal period); measured gaps/jitter and actual delivered rate shall be reported before collection.
\item[NFR-02] Existing model input shall be 200 samples at 100 Hz with six channels unless a new model contract is evaluated.
\item[NFR-03] The archived model-size design budget shall be 120 KB; full firmware flash and RAM shall be measured separately.
\item[NFR-04] Sensitivity and specificity targets of at least 90\% shall be treated as goals, with counts and uncertainty reported; no current result is declared to satisfy both.
\end{description}

